
\needspace{3\baselineskip}
\color{uniblue}{\subsection{Автоматическое регулирование}\label{avtreg}}
\color{black}

\begin{enumerate}[label=\arabic{section}.\arabic{subsection}.\arabic{enumi}, labelsep=4pt, leftmargin=0pt, itemindent=57pt, itemsep=0pt, parsep=5pt] % если что можно крутить itemsep

%%%%%%%%%%%%%%%%%%%%%%%%%%%%%%%%%%%%%%%%%%%%%%%%%%%%%%%%%%%%%%%%%%%%% ОБЩИЕ СВЕДЕНИЯ %%%%%%%%%%%%%%%%%%%%%%%%%%%%%%%%%%%%%%%%%%%%%%%%%%%%%%%%%%%%%%

\item Общие сведения

\begin{enumerate}[label=\arabic{section}.\arabic{subsection}.\arabic{enumi}.\arabic*, labelsep=4pt, leftmargin=0em, itemindent=65pt, parsep=0pt]

\item В режиме автоматического регулирования устройство автоматически формирует команды управления приводом с целью поддержания уровня напряжения в заданном диапазоне.

\item
В таблице \ref{avtreg:tbl1} приведены параметры, необходимые для настройки блока автоматического регулирования.

\small
\begin{longtable}{|>{\centering\arraybackslash}m{5.3cm}|>{\centering\arraybackslash}m{3.3cm}|>{\centering\arraybackslash}m{4.2cm}|>{\centering\arraybackslash}m{1.8cm}|>{\centering\arraybackslash}m{1cm}|}
\caption{Параметры для настройки блока автоматического регулирования\hfill\vspace{-0.5\baselineskip}}\label{avtreg:tbl1}\\ 
\hline
\rowcolor{gray!30}
Параметр (Параметр на ИЧМ) & Условное обозначение на схеме & Значение/ Диапазон & Единица измерения & Шаг \\ 
\hline
\endfirsthead
\caption*{\hspace{3pt}\emph{Продолжение таблицы \ref{avtreg:tbl1}\hfill\vspace{-0.5\baselineskip}}} \\ % сделано по ГОСТ 2.105 п.6.8.7
\hline
\rowcolor{gray!30}
Параметр (Параметр на ИЧМ) & Условное обозначение на схеме & Значение/ Диапазон & Единица измерения & Шаг \\ 
%\hline
\endhead
%\hline
\endfoot
%\hline
\endlastfoot
%==+t1*ATCC1|_TTCC2
\centering Центр зоны нечувствительности (Uпод) & \centering Uпод & \centering 0,85 ... 1,45 & \centering о.е. & \centering \arraybackslash 0,01 \\
\hline
\centering Ширина зоны нечувствительности (dUпод) & \centering dUпод & \centering 0,01 ... 0,20 & \centering о.е. & \centering \arraybackslash 0,01 \\
\hline
\centering Максимальное напряжение выше которого выдача команды «убавить» от автоматики происходит с ускоренной выдержкой времени T3 (U\verb|>>|сраб) & \centering U\verb|>>|сраб & \centering 1,00 ... 2,00 & \centering о.е. & \centering \arraybackslash 0,01 \\
\hline
\centering Активное сопротивление сети 1 секции для линейной компенсации (Rлинии1секции) & \centering Rлинии1секции & \centering 0,00 ... 99999,99 & \centering Ом & \centering \arraybackslash 0,01 \\
\hline
\centering Реактивное сопротивление сети 1 секции для линейной компенсации (Xлинии1секции) & \centering Xлинии1секции & \centering 0,00 ... 99999,99 & \centering Ом & \centering \arraybackslash 0,01 \\
\hline
\centering Активное сопротивление сети 2 секции для линейной компенсации (Rлинии2секции) & \centering Rлинии2секции & \centering 0,00 ... 99999,99 & \centering Ом & \centering \arraybackslash 0,01 \\
\hline
\centering Реактивное сопротивление сети 2 секции для линейной компенсации (Xлинии2секции) & \centering Xлинии2секции & \centering 0,00 ... 99999,99 & \centering Ом & \centering \arraybackslash 0,01 \\
\hline
%===t1
\end{longtable}
\normalsize

\end{enumerate}

%%%%%%%%%%%%%%%%%%%%%%%%%%%%%%%%%%%%%%%%%%%%%%%%%%%%%%%%%%%%%%%%%%%%% Выбор регулируемой секции %%%%%%%%%%%%%%%%%%%%%%%%%%%%%%%%%%%%%%%%%%%%%%%%%%%

\item Выбор регулируемой секции
\begin{enumerate}[label=\arabic{section}.\arabic{subsection}.\arabic{enumi}.\arabic*, labelsep=4pt, leftmargin=0em, itemindent=65pt, parsep=0pt]
\item 
Устройство РПН, установленное на стороне ВН трехобмоточного трансформатора (см. рис. \ref{onsh:img2}) или трансформатора с расщепленной обмоткой  НН (см. рис. \ref{onsh:img1}, б), одновременно воздействует на уровень напряжения на двух подключенных к трансформатору секциях. В данном случае одна из секций определяется как регулируемая, другая -- как контролируемая.
\item
Устройство в автоматическом режиме осуществляет регулирование по уровню напряжения на регулируемой секции. На контролируемой секции, соответственно, проверяется отсутствие условий блокировки автоматического регулирования.
\item
В устройстве предусмотрен следующий механизм определения регулируемой и контролируемой секции. Предусмотрено два режима (определяется состоянием виртуального ключа <<Ввод автоматического выбора регулируемой секции>>):
\begin{itemize}
    \item автоматический;
    \item оперативный.
\end{itemize}

\item
В автоматическом режиме выбор регулируемой секции осуществляется по состоянию блок-контактов вводных выключателей секций и положению виртуального ключа <<Ввод режима приоритета по 2 секции>>. Если подключена только одна из двух секций, она автоматически становится регулируемой.  Если подключены обе секции выбор регулируемой секции осуществляется оперативно виртуальным ключом <<Ввод режима приоритета по 2 секции>>. Подключенная секция, которая не является регулируемой, определяется термином -- контролируемая.В автоматическом режиме выбор регулируемой секции осуществляется по состоянию блок-контактов вводных выключателей секций и положению виртуального ключа <<Ввод режима приоритета по 2 секции>>. Если подключена только одна из двух секций, она автоматически становится регулируемой.  Если подключены обе секции выбор регулируемой секции осуществляется оперативно виртуальным ключом <<Ввод режима приоритета по 2 секции>>. Подключенная секция, которая не является регулируемой, является контролируемой.

\item
В оперативном режиме выбор регулируемой секции осуществляется с помощью виртуального ключа <<Ввод режима приоритета по 2 секции>>. При выведенном состоянии данного виртуального ключа, регулируемой является первая секция, при введенном -- вторая секция. Когда регулируемая секция отключается от трансформатора, то формируется сигнал <<Нет связи с регулируемыми шинами>>. Автоматическое регулирование при этом блокируется.

\item 
Дискретные сигналы <<Секция 1 включена>>, <<Секция 1 отключена>>, <<Секция 2 включена>>, <<Секция 2 отключена>>  могут поступать в логику устройства посредством дискретных входов устройства или GOOSE-сообщения.

\item
Алгоритм выбора регулируемой секции выполнен в соответствии с рисунком \ref{avtreg:img1}.

\vspace{2mm}
\begin{figure}[H]
\centering

\includegraphics[width=1\textwidth,height=1\textheight,keepaspectratio]{\fbpath/502. Авт рег/_latex/img/1.pdf}
\captionof{figure}{Функциональная схема логики выбора регулируемой секции}\label{avtreg:img1}
\end{figure}

\end{enumerate}

%%%%%%%%%%%%%%%%%%%%%%%%%%%%%%%%%%%%%%%%%%%%%%%%%%%%%%%%%%%%%%%%%%%%% Зона нечувствительности %%%%%%%%%%%%%%%%%%%%%%%%%%%%%%%%%%%%%%%%%%%%%%%%%%%

\item Зона нечувствительности
\begin{enumerate}[label=\arabic{section}.\arabic{subsection}.\arabic{enumi}.\arabic*, labelsep=4pt, leftmargin=0em, itemindent=65pt, parsep=0pt]
\item Требуемый диапазон поддержания напряжения на шинах (<<Зона нечувствительности>>) определяется уставками (см. рис. \ref{avtreg:zona}):
\begin{itemize}
    \item центр зоны нечувствительности (или напряжение поддержания) <<Uпод>>;
    \item ширина зоны нечувствительности <<dUпод>>.
\end{itemize}

\vspace{2mm}
\begin{figure}[H]
\centering
\includegraphics[width=0.7\textwidth,height=0.7\textheight,keepaspectratio]{\fbpath/502. Авт рег/_latex/img/zona.png}
\captionof{figure}{Зона нечувствительности}\label{avtreg:zona}
\end{figure}

\item 
Оперативное изменение значения напряжения поддержания в зависимости от суточного графика нагрузки возможно переключением между четырьмя группами уставок. 
\item
Значение напряжения поддержания может меняться автоматически в зависимости от тока нагрузки в режиме встречного регулирования. 
\end{enumerate}

%%%%%%%%%%%%%%%%%%%%%%%%%%%%%%%%%%%%%%%%%%%%%%%%%%%%%%%%%%%%%%%%%%%%% Встречное регулирование %%%%%%%%%%%%%%%%%%%%%%%%%%%%%%%%%%%%%%%%%%%%%%%%%%%

\item Встречное регулирование
\begin{enumerate}[label=\arabic{section}.\arabic{subsection}.\arabic{enumi}.\arabic*, labelsep=4pt, leftmargin=0em, itemindent=65pt, parsep=0pt]
\item
Встречное регулирование –- регулирование напряжения на шинах удаленного потребителя с учетом расчетного падения напряжения в сети в зависимости от тока нагрузки. Для корректного вычисления величины падения напряжения необходимо соблюдать полярность подводимых к устройству цепей переменного тока.
\item 
Падение напряжения в распределительной сети определяется по формуле:

\begin{equation*}
\Delta \dot{U} = \left( \dot{I}{\textsl{св}} - \dot{I}{\textsl{вв}} \right) \cdot \dot{Z}{\textsl{сети}},
\end{equation*}
\begin{eqexpl}[25mm]
    \item{$\dot{Z}_{\textsl{сети}}$} комплексное значение сопротивление сети (линии), определяется по задаваемым параметрам активного и реактивного сопротивления сети, если встречное регулирование не предусматривается, значения данных параметров должны быть равными 0;
    \item{$\dot{I}_{\textsl{вв}}$}комплексное значение тока ввода секции;
    \item{$\dot{I}_{\textsl{св}}$}комплексное значение тока секционного выключателя.
\end{eqexpl}
\item
Значение напряжения на шинах потребителя определяется по формуле:
\begin{equation*}
    {U}_{\textsl{потр}} = \lvert \dot{U}_{\textsl{тек}} - \Delta \dot{U} \rvert,
\end{equation*}

\begin{eqexpl}[25mm]
    \item{$\dot{U}_{\textsl{тек}}$} комплексное значение напряжения на шинах регулируемой секции.
\end{eqexpl}

\end{enumerate}

%%%%%%%%%%%%%%%%%%%%%%%%%%%%%%%%%%%%%%%%%%%%%%%%%%%%%%%%%%%%%%%%%%%%% Процесс автоматического регулирования  %%%%%%%%%%%%%%%%%%%%%%%%%%%%%%%%%%%%%%%%%%%%%%%%%%%

\item Процесс автоматического регулирования 
\begin{enumerate}[label=\arabic{section}.\arabic{subsection}.\arabic{enumi}.\arabic*, labelsep=4pt, leftmargin=0em, itemindent=65pt, parsep=0pt]
\item
В автоматическом режиме устройство непрерывно отслеживает уровень напряжения на регулируемой секции. В устройстве предусмотрены специальные ИО, которые при выходе напряжения за пределы зоны нечувствительности формируют следующие сигналы:
\begin{itemize}
    \item напряжение ниже границы зоны нечувствительности;
    \item напряжение выше границы зоны нечувствительности;
    \item напряжение значительно выше границы зоны нечувствительности.
\end{itemize}
\item
В зависимости от возникающего режима и при отсутствии блокирующих сигналов логика формирования команд в автоматическом режиме регулирования с определенной задержкой формирует сигналы управления приводом.

\end{enumerate}

\end{enumerate}